\documentclass[12pt,a4paper]{report}
\usepackage[ngerman]{babel}
\usepackage{amsmath}
\usepackage{graphicx}

\begin{document}

% ---------- Titelseite ----------
\begin{titlepage}
    \centering
    \vspace*{2cm}
    {\Huge\bfseries Titel der Arbeit\par}
    \vspace{1cm}
    {\Large Untertitel\par}
    \vspace{2cm}
    {\large Vorname Nachname\par}
    \vspace{0.5cm}
    {\large Matrikelnummer: 000000\par}
    \vfill
    {\large Institution\par}
    {\large \today\par}
\end{titlepage}

% ---------- Verzeichnisse ----------
\tableofcontents

% ---------- Kapitel ----------
\chapter{Einleitung}
\label{cha:einleitung}
Hier beginnt die Einleitung. Auf Kapitel~\ref{cha:hauptteil} wird später
verwiesen, und \cite{knuth1984} ist die Standardreferenz.

\chapter{Hauptteil}
\label{cha:hauptteil}
Text des Hauptteils. Eine abgesetzte Formel:
\[
    a^2 + b^2 = c^2 .
\]

\chapter{Fazit}
Zusammenfassung der Ergebnisse.

% ---------- Literatur (offline, ohne biber/bibtex) ----------
\begin{thebibliography}{9}
    \bibitem{knuth1984}
    Donald E.~Knuth. \emph{The \TeX book}. Addison-Wesley, 1984.
\end{thebibliography}

\end{document}
